\documentclass[12pt]{article}
\usepackage[utf8]{inputenc}
\usepackage[T1]{fontenc}
\usepackage{hyperref}
\title{Amerax Language Resources for K-12 Curriculum Development}
\author{Compiled by Claude Sonnet 5.5 (claude-sonnet-5-5) for Andr\'as Kornai}
\date{}
\begin{document}
\maketitle

\section{Introduction}
Amerax (aex) is not a living language that a school curriculum could use. Glottolog files it under ``Bookkeeping'' and records that the ISO 639-3 code was retired and merged into English in 2008. The variety is described as probably English with possible Arabic influence. A German Wikipedia mirror calls it a prison language.

\section{Primary Educational Resources}
None exist.

\section{Linguistic and Reference Materials}
\begin{itemize}
\item Glottolog record: \url{https://glottolog.org/resource/languoid/id/amer1256}
\end{itemize}

\section{Religious and Community Resources}
None found.

\section{Digital and Audio Resources}
None found.

\section{Curriculum Development Context}
No curriculum can be developed for this code. Speakers of the underlying variety would be served by English-language resources.
\end{document}